% =============================================================================
%  Chapter 7 — Element Architecture: Design Patterns and Extensibility
% =============================================================================

\chapter{Element Architecture: Design Patterns and Extensibility}
\label{ch:architecture}

This chapter documents the layered element architecture of \texttt{fall\_n},
focusing on two central abstractions --- \texttt{ElementGeometry} and
\texttt{FEM\_Element} --- their interaction with policy-based design, and the
resulting customisation points.  We compare the approach against the
architectures of OpenSees~\cite{OpenSees} and Kratos
Multiphysics~\cite{Kratos}, highlighting trade-offs in extensibility, compile
time, and runtime efficiency.


% ═══════════════════════════════════════════════════════════════════════════
\section{Guiding Principles}
\label{sec:principles}

\texttt{fall\_n} is designed around three principles that permeate the element
subsystem:

\begin{enumerate}[nosep]
  \item \textbf{Zero-cost abstractions.}  Template instantiation resolves
        element type, quadrature rule, kinematic formulation, and constitutive
        model at compile time.  No virtual dispatch in the inner quadrature
        loop.  The compiler sees the full call chain from Gauss-point
        integration down to the shape function evaluation, enabling inlining
        and SIMD vectorisation.

  \item \textbf{Separation of geometry from physics.}  The geometric
        interpolation (shape functions, Jacobian, quadrature) is decoupled
        from the physical element (strain--displacement operator, constitutive
        law, DOF assembly).  The same \texttt{ElementGeometry<3>} serves both
        a \texttt{ContinuumElement} (3D solid) and a \texttt{SurfaceLoad}
        (boundary traction).

  \item \textbf{Type erasure at system boundaries only.}  Within a single
        element type, everything is monomorphic and stack-allocated.  Virtual
        dispatch is introduced only where heterogeneity is required: storing
        mixed element types in a single container (\texttt{FEM\_Element}) or
        interfacing geometry with the mesh manager
        (\texttt{ElementGeometry::pimpl\_}).
\end{enumerate}

These principles are realised through two complementary patterns:
\emph{policy-based design} (compile-time customisation) and
\emph{external polymorphism} (runtime type erasure).


% ═══════════════════════════════════════════════════════════════════════════
\section{Layer Decomposition}
\label{sec:layers}

The element subsystem is organised in four layers, each with a clear
responsibility boundary:

\begin{center}
\begin{tikzpicture}[
    layer/.style={draw, rounded corners=3pt, minimum width=13cm,
                  minimum height=1.1cm, font=\small},
    arr/.style={-{Stealth[length=5pt]}, thick}
  ]
  \node[layer, fill=blue!8]   (L4) at (0, 4.5)
    {\textbf{Layer 4 — Analysis} : \texttt{Model}, \texttt{LinearAnalysis},
     \texttt{NLAnalysis}, \texttt{DynamicAnalysis}};
  \node[layer, fill=green!8]  (L3) at (0, 3.0)
    {\textbf{Layer 3 — Physical Element} : \texttt{ContinuumElement},
     \texttt{BeamElement}, \texttt{SurfaceLoad}};
  \node[layer, fill=orange!8] (L2) at (0, 1.5)
    {\textbf{Layer 2 — Element Geometry} : \texttt{ElementGeometry<dim>}
     (type-erased handle)};
  \node[layer, fill=red!8]    (L1) at (0, 0.0)
    {\textbf{Layer 1 — Concrete Geometry} : \texttt{LagrangeElement<Dim,N...>},
     \texttt{SimplexElement<Dim,Top,Ord>}};

  \draw[arr] (L1) -- (L2);
  \draw[arr] (L2) -- (L3);
  \draw[arr] (L3) -- (L4);
\end{tikzpicture}
\end{center}

\subsection{Layer~1: Concrete Geometry Types}

These are pure value types with \emph{no virtual functions}.  They encode
the reference element, shape functions, and Jacobian evaluation as
compile-time-resolved operations.

\paragraph{LagrangeElement\texttt{<Dim, N...>}.}
A tensor-product Lagrange element parameterised by the embedding dimension
\texttt{Dim} and per-axis node counts \texttt{N...} (variadic).

\begin{center}
\small
\begin{tabular}{lcl}
  \toprule
  Instantiation & Topology & Gmsh type \\
  \midrule
  \texttt{LagrangeElement<3, 2, 2, 2>} & HEX8  & 5 \\
  \texttt{LagrangeElement<3, 3, 3, 3>} & HEX27 & 12 \\
  \texttt{LagrangeElement<3, 2, 2>}    & QUAD4 (surface) & 3 \\
  \texttt{LagrangeElement<3, 3, 3>}    & QUAD9 (surface) & 10 \\
  \texttt{LagrangeElement<3, 2>}       & LINE2 (edge)    & 1 \\
  \bottomrule
\end{tabular}
\end{center}

Key design features:
\begin{itemize}[nosep]
  \item The reference cell (\texttt{geometry::cell::LagrangianCell<N...>}) is a
        \texttt{constexpr} object holding the nodal coordinates, the
        tensor-product basis, and the subentity topology --- all computed at
        compile time.
  \item The topological dimension is deduced from the parameter pack:
        $\texttt{sizeof...(N)}$.  Thus a surface element is simply a
        \texttt{LagrangeElement<3, 3, 3>} --- same code path, fewer
        template parameters.
  \item Shape functions are products of 1D Lagrange polynomials, evaluated by
        index decomposition.  No lookup table; the compiler folds constant
        expressions.
  \item The \texttt{local\_index\_[]} array enables node reordering (e.g., Gmsh
        to internal convention) without modifying the connectivity array.
\end{itemize}

\paragraph{SimplexElement\texttt{<Dim, TopDim, Order>}.}
A simplex element parameterised by embedding dimension, topological dimension,
and polynomial order.  It mirrors the \texttt{LagrangeElement} public
interface but derives its reference cell from
\texttt{geometry::simplex::SimplexCell<TopDim, Order>}, which stores the
barycentric-based Lagrange shape functions for simplicial geometry.

Both families share a structurally identical interface (node access,
\texttt{H()}, \texttt{dH\_dx()}, \texttt{evaluate\_jacobian()},
\texttt{map\_local\_point()}) but they are \emph{unrelated types} ---
there is no inheritance hierarchy linking them.  Uniformity is achieved
through duck-typing in the next layer.


\subsection{Layer~2: ElementGeometry (Type-Erased Handle)}
\label{sec:elem_geom}

\texttt{ElementGeometry<dim>} is a type-erased value-semantic wrapper that
stores \emph{any} concrete geometry type together with its quadrature rule
behind a single, uniform interface.  It is the pivot of the entire element
subsystem.

\paragraph{Pattern: External Polymorphism.}
The design follows the \emph{External Polymorphism} pattern
(Cleary et al.\ 1997, popularised by Sean Parent's ``Inheritance Is the Base
Class of Evil'' talk):

\begin{enumerate}[nosep]
  \item \textbf{Concept} (\texttt{impl::ElementGeometryConcept<dim>}) ---
        abstract base with pure virtual methods: \texttt{H()},
        \texttt{dH\_dx()}, \texttt{evaluate\_jacobian()},
        \texttt{differential\_measure()}, \texttt{num\_integration\_points()},
        quadrature accessors, and subentity topology.

  \item \textbf{Model} (\texttt{impl::OwningModel\_ElementGeometry%
        <ElementType, IntegrationStrategy>}) ---
        template class deriving from the concept, \emph{owning} a concrete
        \texttt{ElementType} and \texttt{IntegrationStrategy} by value.
        Each virtual override delegates to the stored objects.

  \item \textbf{Handle} (\texttt{ElementGeometry<dim>}) ---
        holds a \texttt{std::unique\_ptr<ElementGeometryConcept<dim>>}
        (\texttt{pimpl\_}).  Provides value semantics via \texttt{clone()},
        move defaulted.
\end{enumerate}

\begin{center}
\begin{tikzpicture}[
    box/.style={draw, rounded corners=2pt, minimum width=4cm,
                minimum height=0.9cm, font=\small\ttfamily},
    arr/.style={-{Stealth[length=4pt]}, thick},
    impl/.style={-{Stealth[length=4pt]}, thick, dashed}
  ]
  \node[box, fill=orange!10] (handle) at (0, 3)
    {ElementGeometry<dim>};
  \node[box, fill=gray!10]   (concept) at (0, 1.5)
    {ElementGeometryConcept<dim>};
  \node[box, fill=green!10]  (model1) at (-3.5, 0)
    {OwningModel<Lag,GL>};
  \node[box, fill=green!10]  (model2) at (3.5, 0)
    {OwningModel<Simp,CP>};

  \draw[impl] (handle) -- node[right, font=\scriptsize]
    {unique\_ptr} (concept);
  \draw[arr]  (model1) -- (concept);
  \draw[arr]  (model2) -- (concept);

  \node[font=\scriptsize, below=2pt of model1]
    {LagrangeElement + GaussLegendre};
  \node[font=\scriptsize, below=2pt of model2]
    {SimplexElement + ConicalProduct};
\end{tikzpicture}
\end{center}

\paragraph{Why External Polymorphism and not classical inheritance?}

\begin{itemize}[nosep]
  \item \texttt{LagrangeElement} and \texttt{SimplexElement} remain pure value
        types --- no vtable pointer, no heap allocation, fully
        \texttt{constexpr}-eligible.  They can be used independently of the
        type-erasure layer (e.g., in constexpr unit tests).
  \item The virtual dispatch lives in the \emph{wrapper}, not in the element
        itself.  Concrete element code retains full optimisation potential.
  \item New element families (e.g., NURBS, serendipity, spectral) are added by
        writing a new value type with the required interface --- no base class
        to inherit from, no code changes in existing elements.
\end{itemize}

\paragraph{Differential measure dispatch.}
The \texttt{differential\_measure()} method in \texttt{OwningModel} uses
\texttt{if constexpr} on the compile-time-known topological dimension:
\begin{lstlisting}[language=C++]
if constexpr (topological_dim == dim)
    return std::abs(J.determinant());        // volume
else if constexpr (topological_dim == 1)
    return J.col(0).norm();                  // arc length
else if constexpr (topological_dim == 2 && dim == 3)
    return J.col(0).cross(J.col(1)).norm();  // surface area
\end{lstlisting}
Despite crossing the virtual boundary, the Jacobian matrix itself is
stack-allocated (Eigen fixed-size \texttt{Matrix<double, dim, Dynamic, \ldots,
dim, dim>}), so no heap allocation occurs inside the hot loop.

\paragraph{Integration via template callback.}
\texttt{ElementGeometry} exposes a non-virtual template method
\texttt{integrate(f)} that accepts any callable \texttt{f(LocalCoordView)}
and loops over quadrature points:
\begin{lstlisting}[language=C++]
template<std::invocable<LocalCoordView> F>
auto integrate(F&& f) const {
    // ... sum w_i * |J(xi_i)| * f(xi_i) ...
}
\end{lstlisting}
Because \texttt{f} is a template parameter (not \texttt{std::function}), the
compiler inlines the integrand.  Only the quadrature point data (coordinates,
weights) and the Jacobian evaluation pass through the virtual boundary; the
user's element-level computation stays in the monomorphic world.


\subsection{Layer~3: Physical Elements}
\label{sec:physical_elements}

Physical elements combine an \texttt{ElementGeometry} with material data and
DOF assembly logic.  Each physical element type is a class template
parameterised by a \emph{MaterialPolicy} and a \emph{KinematicPolicy}.

\paragraph{ContinuumElement\texttt{<MaterialPolicy, ndof, KinematicPolicy>}.}
The workhorse for 1D/2D/3D solid mechanics:
\begin{itemize}[nosep]
  \item Holds a non-owning pointer to an \texttt{ElementGeometry<dim>} (owned
        by the \texttt{Domain}).
  \item Owns the \texttt{MaterialPoint} array (one per Gauss point), each
        holding a \texttt{Material<MaterialPolicy>} with its constitutive
        strategy.
  \item Provides \texttt{inject\_K()}, \texttt{compute\_internal\_forces()},
        \texttt{inject\_tangent\_stiffness()}, \texttt{commit\_material\_state()}
        --- the full \texttt{FiniteElement} concept interface.
  \item The strain--displacement operator $\mathbf{B}$ is delegated to the
        \texttt{KinematicPolicy} (e.g., \texttt{SmallStrain} or
        \texttt{TotalLagrangian}), resolved at compile time.
\end{itemize}

\paragraph{BeamElement\texttt{<BeamPolicy, Dim, AsmPolicy>}.}
Timoshenko beam with section integration:
\begin{itemize}[nosep]
  \item Two-node, local-frame kinematics with a rotation matrix $\mathbf{R}$.
  \item Integration along the beam axis uses
        \texttt{MaterialSection<BeamPolicy>} objects instead of
        \texttt{MaterialPoint}.
  \item The \texttt{AsmPolicy} template parameter selects between direct
        assembly (\texttt{DirectAssembly}) and static condensation
        (\texttt{CondensedAssembly}) --- a compile-time strategy switch.
\end{itemize}

\paragraph{SurfaceLoad.}
Not a ``physical element'' in the structural sense, but uses the
\texttt{ElementGeometry} quadrature machinery to compute consistent nodal
forces from a distributed traction on a boundary face.


\subsection{Layer~4: Analysis Level}
\label{sec:analysis_layer}

\texttt{Model<MaterialPolicy, KinematicPolicy, ndofs, ElemPolicy>} owns the
element container, the \texttt{Domain}, and the PETSc objects
(\texttt{PetscSection}, matrices, vectors).

The \texttt{ElemPolicy} template parameter controls element storage
(\cref{sec:elem_policy}).  For homogeneous meshes, the container is a raw
\texttt{std::vector<ContinuumElement<\ldots>>} with zero overhead.  For
mixed meshes, it is \texttt{std::vector<FEM\_Element>} with one virtual
dispatch per element call.


% ═══════════════════════════════════════════════════════════════════════════
\section{FEM\_Element: Top-Level Type Erasure}
\label{sec:fem_element}

\texttt{FEM\_Element} applies the same External Polymorphism pattern as
\texttt{ElementGeometry}, but at the assembly level:

\begin{center}
\begin{tikzpicture}[
    box/.style={draw, rounded corners=2pt, minimum width=4.2cm,
                minimum height=0.9cm, font=\small\ttfamily},
    arr/.style={-{Stealth[length=4pt]}, thick},
    impl/.style={-{Stealth[length=4pt]}, thick, dashed}
  ]
  \node[box, fill=orange!10] (handle) at (0, 3)
    {FEM\_Element};
  \node[box, fill=gray!10]   (concept) at (0, 1.5)
    {Concept (inner struct)};
  \node[box, fill=green!10]  (model1) at (-3.5, 0)
    {Model<ContinuumElem>};
  \node[box, fill=green!10]  (model2) at (3.5, 0)
    {Model<BeamElement>};

  \draw[impl] (handle) -- node[right, font=\scriptsize]
    {unique\_ptr} (concept);
  \draw[arr]  (model1) -- (concept);
  \draw[arr]  (model2) -- (concept);
\end{tikzpicture}
\end{center}

Key design decisions:

\begin{itemize}[nosep]
  \item \textbf{Value semantics.} Copy via \texttt{clone()}, move defaulted.
        No shared ownership, no reference counting.
  \item \textbf{Self-wrap guard.} A \texttt{requires} clause prevents wrapping
        a \texttt{FEM\_Element} inside another \texttt{FEM\_Element} (double
        indirection).
  \item \textbf{Recursive concept satisfaction.} \texttt{FEM\_Element} itself
        satisfies the \texttt{FiniteElement} concept, verified by a
        \texttt{static\_assert} at the bottom of the header.  This means it
        can be nested inside other type-erased wrappers without losing the
        concept constraint.
  \item \textbf{Optional interface methods.} Mass matrix operations
        (\texttt{density()}, \texttt{inject\_mass()}) use \texttt{if constexpr}
        inside \texttt{Model<T>} to detect whether the concrete element
        supports them, defaulting to a no-op otherwise.  This avoids bloating
        the concept with operations not all elements provide.
\end{itemize}

There is also a \texttt{StructuralElement} wrapper, identical in pattern,
that provides type safety for structural-only models (beams, shells).  Both
derive from \emph{no common base class} --- they are independent type-erasure
handles, unified only by the \texttt{FiniteElement} concept.


% ═══════════════════════════════════════════════════════════════════════════
\section{Policy-Based Design: Customisation Points}
\label{sec:policies}

\texttt{fall\_n} uses policies --- template parameters that select
compile-time strategies --- at multiple levels.  Each policy is a
\emph{customisation point}: a seam where the user can inject different
behaviour without modifying existing code.

\subsection{MaterialPolicy}
\label{sec:mat_policy}

A \emph{trait bag} that binds kinematic and conjugate types:

\begin{center}
\small
\begin{tabular}{lllc}
  \toprule
  Policy & \texttt{StrainT} & \texttt{StressT} & \texttt{dim} \\
  \midrule
  \texttt{ThreeDimensionalMaterial} & \texttt{Strain<6>} & \texttt{Stress<6>} & 3 \\
  \texttt{PlaneMaterial}            & \texttt{Strain<3>} & \texttt{Stress<3>} & 2 \\
  \texttt{UniaxialMaterial}         & \texttt{Strain<1>} & \texttt{Stress<1>} & 1 \\
  \texttt{TimoshenkoBeam3D}         & \texttt{BeamStrain<6,3>} & \texttt{BeamForces<6>} & 3 \\
  \texttt{TimoshenkoBeam2D}         & \texttt{BeamStrain<3,2>} & \texttt{BeamForces<3>} & 2 \\
  \bottomrule
\end{tabular}
\end{center}

All generic code (\texttt{Material<P>}, \texttt{MaterialPoint<P>},
\texttt{MaterialState<P>}, \texttt{ElasticRelation<P>}) is parameterised
on a single \texttt{MaterialPolicy} and extracts its types via
\texttt{P::StrainT}, \texttt{P::StressT}, etc.

\paragraph{Extensibility.}
Adding a shell element requires only:
\begin{enumerate}[nosep]
  \item Define \texttt{ShellMaterial<N, 3>} with appropriate
        \texttt{ShellGeneralizedStrain<N,3>} and
        \texttt{ShellResultants<N>}.
  \item Write a constitutive model conforming to
        \texttt{ConstitutiveRelation<ShellMaterial<N,3>>}.
  \item Write the element class (\texttt{ShellElement}) satisfying
        \texttt{FiniteElement}.
\end{enumerate}
No existing code changes; the policy mechanism propagates types automatically.


\subsection{KinematicPolicy}
\label{sec:kin_policy}

Selects the strain--displacement relationship at compile time:

\begin{center}
\small
\begin{tabular}{lp{5cm}cc}
  \toprule
  Policy & Description & Linear? & $K_\sigma$? \\
  \midrule
  \texttt{SmallStrain}       & $\bvarepsilon = \nabla^s \mathbf{u}$ & yes & no \\
  \texttt{TotalLagrangian}   & $\mathbf{E} = \frac{1}{2}(\bF^T\bF - \bI)$ & no & yes \\
  \texttt{UpdatedLagrangian} & Cauchy strain in current config. & no & yes \\
  \texttt{Corotational}      & Large rotation, small strain & no & yes \\
  \bottomrule
\end{tabular}
\end{center}

Each policy is a struct with static methods:
\begin{lstlisting}[language=C++]
struct SmallStrain {
    static constexpr bool is_geometrically_linear    = true;
    static constexpr bool needs_geometric_stiffness  = false;
    
    template <std::size_t dim>
    static GPKinematics<dim> evaluate(
        ElementGeometry<dim>* geo, std::size_t num_nodes,
        std::size_t ndof, const Array& Xi,
        const Eigen::VectorXd& u_e);

    template <std::size_t dim>
    static auto compute_B(ElementGeometry<dim>* geo, ...);
};
\end{lstlisting}

\texttt{ContinuumElement} calls \texttt{KinematicPolicy::evaluate<dim>()} at
each quadrature point.  The boolean traits steer \texttt{if constexpr}
branches inside the element (e.g., add $K_\sigma$ only when
\texttt{needs\_geometric\_stiffness == true}).

\paragraph{Extensibility.}
A new formulation (e.g., logarithmic strain) is added by defining a new
struct satisfying \texttt{KinematicPolicyConcept} with the required static
methods and boolean traits.  No changes to \texttt{ContinuumElement} or any
analysis class.


\subsection{ElementPolicy (Storage Strategy)}
\label{sec:elem_policy}

Controls how \texttt{Model} stores its element collection:

\begin{center}
\small
\begin{tabular}{lp{4.5cm}cc}
  \toprule
  Policy & Container & Homogeneous & Overhead \\
  \midrule
  \texttt{SingleElementPolicy<E>} & \texttt{std::vector<E>}
    & yes & none \\
  \texttt{MultiElementPolicy}     & \texttt{std::vector<FEM\_Element>}
    & no  & 1 vtable + 1 ptr chase \\
  \bottomrule
\end{tabular}
\end{center}

The concept \texttt{ElementPolicyLike} enforces the policy interface:
\begin{lstlisting}[language=C++]
template <typename P>
concept ElementPolicyLike = requires {
    typename P::element_type;
    typename P::container_type;
    { P::is_homogeneous } -> std::convertible_to<bool>;
} && FiniteElement<typename P::element_type>;
\end{lstlisting}

\paragraph{Default path.}
$\texttt{Model<ThreeDimensionalMaterial>}$ defaults to:\\
$\texttt{SingleElementPolicy<ContinuumElement<ThreeDimensionalMaterial, 3,}$
$\texttt{SmallStrain>>}$,\\
yielding a contiguous \texttt{std::vector} of monomorphic elements --- optimal
for cache locality and auto-vectorisation.


\subsection{AssemblyPolicy (Beams)}

Beam elements accept an \texttt{AsmPolicy} parameter:
\begin{itemize}[nosep]
  \item \texttt{DirectAssembly} --- full stiffness matrix, no condensation.
  \item \texttt{CondensedAssembly} --- static condensation of internal DOFs
        before injection into PETSc.
\end{itemize}
Both are stateless tag types; the choice is resolved at compile time via
\texttt{if constexpr} inside \texttt{BeamElement::inject\_K()}.

\subsection{ConstitutiveStrategy (Material Update)}

The material update mechanism (\texttt{Material<P>}) delegates to a
type-erased strategy object (\texttt{IntegrationStrategy<P>}), enabling
runtime polymorphism over the constitutive response:
\begin{itemize}[nosep]
  \item \texttt{ElasticRelation<P>} --- linear elastic.
  \item \texttt{HyperelasticRelation<P>} --- Neo-Hookean, SVK, etc.
  \item \texttt{J2Plasticity<P>} --- return mapping with isotropic hardening.
\end{itemize}
This is the \emph{only} place in the element quadrature loop where runtime
polymorphism is used, because the constitutive model is genuinely a runtime
choice (same mesh, different materials per physical group).


% ═══════════════════════════════════════════════════════════════════════════
\section{Compile-Time Cost Analysis}
\label{sec:compile_time}

Policy-based design with heavy template usage has a well-known compile-time
cost.  \texttt{fall\_n} mitigates this through several strategies:

\begin{enumerate}[nosep]
  \item \textbf{Header-only by necessity, not by choice.}  The library is
        header-only because C++ templates must be visible at the point of
        instantiation.  However, explicit template instantiation could be
        added in \texttt{.cpp} files for the most common configurations to
        reduce redundant codegen across translation units.

  \item \textbf{Type erasure as a firewall.}  \texttt{ElementGeometry<dim>}
        and \texttt{FEM\_Element} act as \emph{template firewalls}: code that
        consumes an \texttt{ElementGeometry<3>} does not need to see the
        headers of \texttt{LagrangeElement} or \texttt{SimplexElement}.
        This significantly reduces transitive include graphs.

  \item \textbf{Consteval reference cells.}  The \texttt{LagrangianCell<N...>}
        and \texttt{SimplexCell<TopDim, Order>} objects with their basis
        functions and topological tables are \texttt{constexpr}/\texttt{consteval},
        evaluated by the compiler once.  They produce no runtime code unless
        actually called; the compiler may fold constant subexpressions entirely.

  \item \textbf{Concept constraints as documentation.}  The
        \texttt{FiniteElement}, \texttt{KinematicPolicyConcept},
        \texttt{ElementPolicyLike} concepts serve as lightweight checks that
        produce clear error messages at instantiation time, avoiding the
        traditional deep-template-error-message problem.
\end{enumerate}

\paragraph{Measured impact.}
A full rebuild of \texttt{fall\_n} (main + 12 test targets) takes approximately
105~seconds on a single core with GCC~11 and debug optimisation (\texttt{-g}).
The dominant cost is Eigen's template instantiation inside
\texttt{ContinuumElement} and \texttt{BeamElement}; the element geometry layer
adds minimal overhead because the type-erasure boundary limits template
propagation.


% ═══════════════════════════════════════════════════════════════════════════
\section{Comparison with OpenSees and Kratos}
\label{sec:comparison}

\subsection{Overview}

\begin{center}
\small
\renewcommand{\arraystretch}{1.2}
\begin{tabular}{p{3.8cm}p{3.8cm}p{3.8cm}p{3.8cm}}
  \toprule
  Aspect & \textbf{OpenSees} & \textbf{Kratos} & \textbf{fall\_n} \\
  \midrule
  Language &
    C++ (C++11/14) &
    C++ (C++17) + Python &
    C++ (C++20/23) \\
  Polymorphism model &
    Classical inheritance &
    Classical inheritance + Python bridge &
    Policies + External Polymorphism \\
  Element base &
    \texttt{Element} (virtual) &
    \texttt{Element} (virtual) &
    \texttt{FiniteElement} (concept) \\
  Geometry base &
    \texttt{CrdTransf} (virtual) &
    \texttt{Geometry<TPoint>} (virtual) &
    \texttt{ElementGeometry<dim>} (type-erased) \\
  Material base &
    \texttt{UniaxialMaterial} / \texttt{NDMaterial} (virtual) &
    \texttt{ConstitutiveLaw} (virtual) &
    \texttt{Material<Policy>} (policy + type-erased strategy) \\
  Dispatch cost &
    1--3 vtable calls per GP &
    1--2 vtable calls per GP &
    0 (homogeneous) or 1 (heterogeneous) \\
  \bottomrule
\end{tabular}
\end{center}

\subsection{OpenSees}

OpenSees~\cite{OpenSees} uses a deep inheritance hierarchy:
$\texttt{DomainComponent} \to \texttt{Element} \to$ concrete types.
Materials form an even deeper tree:
$\texttt{Material} \to \texttt{UniaxialMaterial} \to$ specific laws.

\paragraph{Strengths.}
\begin{itemize}[nosep]
  \item Extremely mature (20+ years), battle-tested in earthquake engineering.
  \item Simple runtime scripting via Tcl (later Python) --- users define
        models without recompiling.
  \item Large element library (200+ elements).
\end{itemize}

\paragraph{Weaknesses (from an architecture perspective).}
\begin{itemize}[nosep]
  \item \textbf{Monolithic virtual hierarchy.}  All elements inherit from
        \texttt{Element}, which carries 30+ virtual methods including
        \texttt{getResistingForce()}, \texttt{getTangentStiff()}, etc.
        Adding a new element requires implementing all of them (or relying
        on default no-ops that may silently produce incorrect results).
  \item \textbf{No separation of geometry and physics.}  A \texttt{Brick20N}
        element hard-codes its shape functions, Jacobian, AND constitutive
        integration in a single class.  Changing the quadrature rule requires
        modifying the element source.
  \item \textbf{Dynamic dispatch in the inner loop.}  Every call to
        \texttt{material->getStress()} and \texttt{material->getTangent()}
        at each Gauss point goes through virtual dispatch.  While the cost
        is small relative to the linear algebra, it inhibits inlining and
        auto-vectorisation across the material boundary.
  \item \textbf{No concept or constraint checking.}  Adding a new element
        with a typo in a virtual override signature compiles silently and
        produces incorrect results at runtime.
\end{itemize}

\paragraph{Contrast with fall\_n.}
\texttt{fall\_n} eliminates the monolithic base class by defining
\texttt{FiniteElement} as a C++20 concept that enforces the interface at
compile time.  The element is free to organise its internals however it
pleases; the concept checks only what the assembler needs.  Geometry is
factored out into \texttt{ElementGeometry}, so the same mesh can serve
continuum, surface-load, and post-processing operations.


\subsection{Kratos Multiphysics}

Kratos~\cite{Kratos} uses a more modular architecture than OpenSees, with
explicit separation of geometry (\texttt{Geometry<TPoint>}), element
(\texttt{Element}), condition (\texttt{Condition}), constitutive law
(\texttt{ConstitutiveLaw}), and process (\texttt{Process}).

\paragraph{Strengths.}
\begin{itemize}[nosep]
  \item \textbf{Multi-physics by design.}  The \texttt{ModelPart} /
        \texttt{ProcessInfo} / \texttt{Variable} registry enables coupling of
        disparate physics (structural, fluid, thermal) sharing the same mesh.
  \item \textbf{Python-driven workflow.}  Element selection, boundary
        conditions, solver parameters, and post-processing are all scriptable
        in Python without recompilation.
  \item \textbf{Geometry separated from element.}  \texttt{Geometry<TPoint>}
        handles shape functions and Jacobian; \texttt{Element} handles physics.
        This is conceptually analogous to \texttt{fall\_n}'s two-layer split.
\end{itemize}

\paragraph{Weaknesses (from an architecture perspective).}
\begin{itemize}[nosep]
  \item \textbf{Runtime overhead of the variable system.}
        Each nodal value is stored in a hash-map-like container
        (\texttt{Variable<T>} keys), requiring a lookup per node per variable
        access.  This sacrifices cache locality and prevents compile-time
        resolution of DOF layouts.
  \item \textbf{Virtual geometry hierarchy.}  \texttt{Geometry<TPoint>} is an
        inheritance tree: $\texttt{Geometry} \to \texttt{Triangle2D3}$, etc.
        Shape functions are declared \texttt{virtual}, adding dispatch cost at
        each Gauss point.  There is no mechanism to ``inline'' the geometry into
        the element at compile time.
  \item \textbf{Heavyweight element interface.}  \texttt{Element} carries
        methods for \texttt{CalculateLocalSystem},
        \texttt{CalculateMassMatrix}, \texttt{CalculateDampingMatrix},
        \texttt{GetDofList}, etc.  Many elements implement subsets, leaving
        the rest as no-ops or throwing exceptions --- a code-smell of
        interface bloat.
  \item \textbf{No compile-time guarantees.}  Registering a new element
        involves adding entries to a factory map at runtime.  Errors in the
        element's interface (wrong return type, missing method) are caught
        only at link time or, worse, at runtime.
\end{itemize}

\paragraph{Contrast with fall\_n.}
\texttt{fall\_n}'s policy-based design replaces the Kratos variable registry
with compile-time type parameters: the \texttt{MaterialPolicy} selects the
strain/stress types, the \texttt{KinematicPolicy} selects the B-operator, and
the \texttt{ElementPolicy} selects the storage strategy --- all resolved by
the compiler.  The trade-off is that \texttt{fall\_n} requires recompilation
to change element types, whereas Kratos allows runtime scripting.


\subsection{Summary: Architectural Trade-offs}
\label{sec:tradeoffs}

\begin{center}
\small
\renewcommand{\arraystretch}{1.2}
\begin{tabular}{p{4cm}ccc}
  \toprule
  Trade-off & OpenSees & Kratos & fall\_n \\
  \midrule
  Runtime scripting         & \checkmark\checkmark & \checkmark\checkmark\checkmark & --- \\
  Compile-time safety       & ---  & ---  & \checkmark\checkmark\checkmark \\
  Zero-cost inner loop      & ---  & ---  & \checkmark\checkmark \\
  Geometry--physics split   & ---  & \checkmark\checkmark & \checkmark\checkmark\checkmark \\
  Heterogeneous elements    & \checkmark\checkmark\checkmark & \checkmark\checkmark\checkmark & \checkmark\checkmark \\
  Multi-physics coupling    & \checkmark & \checkmark\checkmark\checkmark & --- \\
  Maturity / element count  & \checkmark\checkmark\checkmark & \checkmark\checkmark\checkmark & \checkmark \\
  \bottomrule
\end{tabular}
\end{center}

\texttt{fall\_n} occupies a different niche: it is \emph{not} a general-purpose
production FE framework, but a research library where the primary concern is
\textbf{architectural clarity and compile-time correctness} over runtime
flexibility.  The policy-based approach is intentionally opinionated: it makes
the ``easy things trivial'' (adding a new constitutive model, kinematic
formulation, or element family within the existing type system) at the cost of
making the ``hard things hard'' (mixing unrelated physics, runtime element
selection from user input).


% ═══════════════════════════════════════════════════════════════════════════
\section{Extensibility Guide: Adding a New Element Family}
\label{sec:extensibility}

This section walks through the concrete steps to add a hypothetical
``Serendipity element'' family, illustrating each customisation point.

\begin{enumerate}
  \item \textbf{Layer~1: Concrete geometry type.}
    Create \texttt{SerendipityElement<Dim, Order>}:
    \begin{itemize}[nosep]
      \item Store node coordinates, define \texttt{H()}, \texttt{dH\_dx()},
            \texttt{evaluate\_jacobian()}, \texttt{map\_local\_point()} with
            the serendipity basis functions.
      \item No base class needed.  The type must be movable and copyable.
    \end{itemize}

  \item \textbf{Layer~1: Integration strategy.}
    Create or reuse a quadrature class (\texttt{GaussLegendre2D<p>},
    \texttt{ReducedIntegration<\ldots>}, etc.) exposing
    \texttt{reference\_integration\_point(i)} and \texttt{weight(i)}.

  \item \textbf{Layer~2: Instantiate through ElementGeometry.}
    No code changes.  The constructor template:
\begin{lstlisting}[language=C++]
ElementGeometry<3> geo(
    SerendipityElement<3, 2>{tag, node_ids},
    GaussLegendre2D<3>{});
\end{lstlisting}
    automatically instantiates an
    \texttt{OwningModel<Serendipity..., GL...>} and stores it behind the
    type-erased handle.

  \item \textbf{Layer~3: Use with ContinuumElement.}
    No code changes.  \texttt{ContinuumElement} receives a pointer to the
    geometry and calls \texttt{H()}, \texttt{dH\_dx()},
    \texttt{differential\_measure()} through the virtual interface.  The
    $\mathbf{B}$ matrix and assembly are generic.

  \item \textbf{Layer~4: Mesh import.}
    Add a case in \texttt{GmshDomainBuilder.hh} for the Gmsh element type
    code, constructing \texttt{ElementGeometry} with the serendipity concrete
    type and the desired integrator.

  \item \textbf{Optional: Type-safe model.}
    If the serendipity element will be used in its own model (no mixing),
    define:
\begin{lstlisting}[language=C++]
using SerendipitySolid =
    ContinuumElement<ThreeDimensionalMaterial, 3,
                     SmallStrain>;
using SerendipityModel =
    Model<ThreeDimensionalMaterial, SmallStrain, 3,
          SingleElementPolicy<SerendipitySolid>>;
\end{lstlisting}
    This compiles to a contiguous vector with zero virtual dispatch.
\end{enumerate}

\paragraph{What does NOT need to change:}
\begin{itemize}[nosep]
  \item \texttt{ElementGeometry.hh} (new type is auto-detected by the
        templated constructor).
  \item \texttt{ContinuumElement.hh} (works through geometry pointer).
  \item \texttt{FEM\_Element.hh} (type erasure is generic).
  \item \texttt{Model.hh} (policy-parameterised).
  \item Any analysis class (\texttt{LinearAnalysis}, \texttt{NLAnalysis}, etc.).
  \item Any material or constitutive model.
\end{itemize}

This is the central extensibility promise: \textbf{new element families are
additive} --- they require only new files, not modifications to existing ones
--- conforming to the Open/Closed Principle.


% ═══════════════════════════════════════════════════════════════════════════
\section{Runtime Efficiency Considerations}
\label{sec:efficiency}

\subsection{Virtual Dispatch Budget}

In the homogeneous path (\texttt{SingleElementPolicy}), the element assembly
loop has \textbf{zero} virtual calls for the physics: all
\texttt{inject\_K()}, \texttt{B()}, \texttt{H()}, \texttt{dH\_dx()} calls are
monomorphic and inlineable.

The only virtual calls per Gauss point are:
\begin{itemize}[nosep]
  \item \texttt{ElementGeometry::evaluate\_jacobian()} --- 1 vcall for the
        Jacobian (through pimpl).
  \item \texttt{ElementGeometry::H()}, \texttt{dH\_dx()} --- $n_\text{nodes}
        \times (1 + \text{dim})$ vcalls for shape functions and derivatives.
  \item \texttt{Material::compute\_response()} --- 1 vcall for the
        constitutive model (through the strategy pointer).
\end{itemize}
This is \emph{identical per-element overhead} to OpenSees or Kratos, but
\texttt{fall\_n} avoids the additional vtable dispatch that those frameworks
incur at the element container level.

\subsection{Cache-Friendly Layout}

\texttt{SingleElementPolicy<E>} stores all elements contiguously in a single
\texttt{std::vector<E>}.  Each element's \texttt{MaterialPoint} array is
heap-allocated (as a \texttt{std::vector<MaterialPointT>}), but its entries
are contiguous for that element.

\texttt{MultiElementPolicy} stores \texttt{FEM\_Element} objects contiguously
(the \texttt{unique\_ptr<Concept>} wrappers are packed), but the pointed-to
\texttt{Model<T>} objects are heap-allocated with potentially scattered
addresses.  This adds one pointer chase per element --- acceptable for
heterogeneous cases where the element count is moderate.

\subsection{DOF Index Caching}

\texttt{ContinuumElement} caches its flat \texttt{dof\_indices\_[]} array
(one entry per element DOF, combining all nodes).  The cache is populated
lazily on first use and reused across \texttt{inject\_K()},
\texttt{compute\_internal\_forces()}, and \texttt{inject\_tangent\_stiffness()}.
This avoids the per-call overhead of traversing the node pointers to
collect DOF indices at every assembly step.
\begin{lstlisting}[language=C++]
void ensure_dof_cache() noexcept {
    if (dofs_cached_) return;
    collect_dof_indices();
}
\end{lstlisting}

\subsection{Stack-Allocated Jacobians}

The Jacobian matrix crossing the virtual boundary is declared as:
\begin{lstlisting}[language=C++]
using JacobianMatrix = Eigen::Matrix<double, dim,
    Eigen::Dynamic, Eigen::ColMajor, dim, dim>;
\end{lstlisting}
The \texttt{MaxRows = dim}, \texttt{MaxCols = dim} template parameters tell
Eigen to use a fixed-size stack buffer (up to $3 \times 3 = 72$ bytes) with
a dynamic column count.  Since \texttt{dim} $\le 3$, no heap allocation
occurs, even though the column count is ``dynamic'' (it varies between volume
and surface elements).

This ensures that the virtual call to \texttt{evaluate\_jacobian()} returns by
value without any \texttt{new}/\texttt{delete}, making the hot path
allocation-free.
